% !TEX program = xelatex
% 城市内涝实时预警系统 —— 多图元组合（单色 / 彩色两版）
\documentclass[10pt]{article}
\usepackage[a4paper,margin=1.5cm]{geometry}
\usepackage{xeCJK}
\setCJKmainfont{Microsoft YaHei}[BoldFont={Microsoft YaHei Bold}]
\setmainfont{Microsoft YaHei}[BoldFont={Microsoft YaHei Bold}]
\xeCJKsetup{CJKecglue={}}
\usepackage{tikz}
\usetikzlibrary{arrows.meta}
\definecolor{monoLine}{HTML}{4D4D4D}
\definecolor{monoHead}{HTML}{D9D9D9}
\definecolor{monoSide}{HTML}{CBCBCB}
\definecolor{navyDeep}{HTML}{2A4370}
\definecolor{navyLight}{HTML}{D6E0EE}
\definecolor{navySide}{HTML}{B9CBE2}
\definecolor{tealDeep}{HTML}{2C736E}
\definecolor{tealLight}{HTML}{D3E7E4}
\definecolor{tealSide}{HTML}{B7D8D5}
\definecolor{warmDeep}{HTML}{B0702F}
\definecolor{warmLight}{HTML}{F4E6D4}
\definecolor{warmSide}{HTML}{E3C7A4}
\setlength{\parindent}{0pt}
\pagestyle{empty}
\begin{document}

\begin{center}
{\large\bfseries 城市内涝实时预警系统技术路线}\\[1pt]
{\footnotesize 同一份几何，两种配色}
\end{center}

\vspace{0pt}
{\footnotesize\textbf{(a) 单色版}　黑白印刷稿用。}

\vspace{0pt}
\begin{center}
\begin{tikzpicture}[x=0.0138889cm,y=-0.0138889cm]
\draw[dashed,dash pattern=on 2.0pt off 1.6pt,line width=0.7pt,draw=monoLine] (40.0,40.0) rectangle (960.0,404.0);
\draw[line width=0.7pt,draw=monoLine] (54.0,390.0) -- (54.0,54.0) -- (86.0,54.0) -- (86.0,210.0);
\draw[line width=0.7pt,fill=white,draw=monoLine] (86.0,210.0) -- (118.0,210.0) -- (118.0,197.0) -- (160.0,222.0) -- (118.0,247.0) -- (118.0,234.0) -- (86.0,234.0);
\draw[line width=0.7pt,draw=monoLine] (86.0,234.0) -- (86.0,390.0) -- (54.0,390.0);
\node[align=center,inner sep=0pt,outer sep=0pt,font=\fontsize{8.6}{12.04}\selectfont,text=monoLine] at (70.0,153.0) {数};
\node[align=center,inner sep=0pt,outer sep=0pt,font=\fontsize{8.6}{12.04}\selectfont,text=monoLine] at (70.0,199.0) {据};
\node[align=center,inner sep=0pt,outer sep=0pt,font=\fontsize{8.6}{12.04}\selectfont,text=monoLine] at (70.0,245.0) {治};
\node[align=center,inner sep=0pt,outer sep=0pt,font=\fontsize{8.6}{12.04}\selectfont,text=monoLine] at (70.0,291.0) {理};
\draw[fill=white,draw=monoLine,line width=0.7pt] (168.0,68.0) -- (168.0,166.0) .. controls (168.0,171.0) and (346.0,171.0) .. (346.0,166.0) -- (346.0,68.0) -- cycle;
\draw[fill=white,draw=monoLine,line width=0.7pt] (257.0,68.0) ellipse (89.0 and 9.0);
\node[align=center,inner sep=0pt,outer sep=0pt,font=\fontsize{8.3}{11.62}\selectfont,text=monoLine] at (257.0,123.0) {雨量遥测库};
\draw[fill=white,draw=monoLine,line width=0.7pt] (368.0,68.0) -- (368.0,166.0) .. controls (368.0,171.0) and (546.0,171.0) .. (546.0,166.0) -- (546.0,68.0) -- cycle;
\draw[fill=white,draw=monoLine,line width=0.7pt] (457.0,68.0) ellipse (89.0 and 9.0);
\node[align=center,inner sep=0pt,outer sep=0pt,font=\fontsize{8.3}{11.62}\selectfont,text=monoLine] at (457.0,123.0) {管网拓扑库};
\draw[fill=white,draw=monoLine,line width=0.7pt] (568.0,68.0) -- (568.0,166.0) .. controls (568.0,171.0) and (746.0,171.0) .. (746.0,166.0) -- (746.0,68.0) -- cycle;
\draw[fill=white,draw=monoLine,line width=0.7pt] (657.0,68.0) ellipse (89.0 and 9.0);
\node[align=center,inner sep=0pt,outer sep=0pt,font=\fontsize{8.3}{11.62}\selectfont,text=monoLine] at (657.0,123.0) {历史内涝库};
\draw[fill=white,draw=monoLine,line width=0.7pt] (768.0,68.0) -- (768.0,166.0) .. controls (768.0,171.0) and (946.0,171.0) .. (946.0,166.0) -- (946.0,68.0) -- cycle;
\draw[fill=white,draw=monoLine,line width=0.7pt] (857.0,68.0) ellipse (89.0 and 9.0);
\node[align=center,inner sep=0pt,outer sep=0pt,font=\fontsize{8.3}{11.62}\selectfont,text=monoLine] at (857.0,123.0) {预案知识库};
\draw[line width=0.7pt,draw=monoLine] (257.0,196.0) -- (857.0,196.0);
\draw[line width=0.7pt,draw=monoLine] (257.0,175.0) -- (257.0,196.0);
\draw[line width=0.7pt,draw=monoLine] (457.0,175.0) -- (457.0,196.0);
\draw[line width=0.7pt,draw=monoLine] (657.0,175.0) -- (657.0,196.0);
\draw[line width=0.7pt,draw=monoLine] (857.0,175.0) -- (857.0,196.0);
\draw[line width=0.7pt,draw=monoLine] (578.0,196.0) -- (578.0,222.0);
\fill[fill=white,draw=monoLine,line width=0.7pt] (257.0,264.0) -- (299.0,222.0) -- (899.0,222.0) -- (857.0,264.0) -- cycle;
\fill[fill=monoSide,draw=monoLine,line width=0.7pt] (857.0,264.0) -- (899.0,222.0) -- (899.0,332.0) -- (857.0,374.0) -- cycle;
\draw[fill=white,draw=monoLine,line width=0.7pt] (257.0,264.0) rectangle (857.0,374.0);
\node[align=center,inner sep=0pt,outer sep=0pt,font=\fontsize{8.8}{12.32}\selectfont,text=monoLine] at (578.0,243.0) {\bfseries 多源数据融合与样本构建};
\node[align=center,inner sep=0pt,outer sep=0pt,font=\fontsize{7.6}{10.64}\selectfont,text=black] at (557.0,310.0) {时钟对齐 · 缺测修复 · 汇水单元编码};
\node[align=center,inner sep=0pt,outer sep=0pt,font=\fontsize{7.6}{10.64}\selectfont,text=black] at (557.0,336.0) {特征构建 · 样本分层 · 训练验证切分};
\draw[dashed,dash pattern=on 2.0pt off 1.6pt,line width=0.7pt,draw=monoLine] (40.0,482.0) rectangle (960.0,764.0);
\draw[line width=0.7pt,draw=monoLine] (54.0,750.0) -- (54.0,496.0) -- (86.0,496.0) -- (86.0,611.0);
\draw[line width=0.7pt,fill=white,draw=monoLine] (86.0,611.0) -- (118.0,611.0) -- (118.0,598.0) -- (160.0,623.0) -- (118.0,648.0) -- (118.0,635.0) -- (86.0,635.0);
\draw[line width=0.7pt,draw=monoLine] (86.0,635.0) -- (86.0,750.0) -- (54.0,750.0);
\node[align=center,inner sep=0pt,outer sep=0pt,font=\fontsize{8.6}{12.04}\selectfont,text=monoLine] at (70.0,554.0) {模};
\node[align=center,inner sep=0pt,outer sep=0pt,font=\fontsize{8.6}{12.04}\selectfont,text=monoLine] at (70.0,600.0) {型};
\node[align=center,inner sep=0pt,outer sep=0pt,font=\fontsize{8.6}{12.04}\selectfont,text=monoLine] at (70.0,646.0) {推};
\node[align=center,inner sep=0pt,outer sep=0pt,font=\fontsize{8.6}{12.04}\selectfont,text=monoLine] at (70.0,692.0) {演};
\draw[draw=monoLine,line width=0.7pt,fill=white,rounded corners=0.1667cm] (180.0,528.0) rectangle (330.0,618.0);
\node[align=center,inner sep=0pt,outer sep=0pt,text=monoLine] at (255.0,573.0) {{\fontsize{8.30}{9.54}\selectfont\bfseries 模型输入}\\[1.20pt]{\fontsize{6.40}{7.68}\selectfont 边界与初始场}};
\draw[draw=monoLine,line width=0.7pt,fill=white,rounded corners=0.1667cm] (180.0,660.0) rectangle (330.0,750.0);
\node[align=center,inner sep=0pt,outer sep=0pt,text=monoLine] at (255.0,705.0) {{\fontsize{8.30}{9.54}\selectfont\bfseries 参数标定}\\[1.20pt]{\fontsize{6.40}{7.68}\selectfont 糙率 · 坡度 · 管径}};
\draw[line width=0.7pt,draw=monoLine,-{Latex[length=1.5mm,width=1.0mm]}] (255,618.0) -- (255,660.0);
\draw[draw=monoLine,line width=0.7pt,fill=white,rounded corners=0.1389cm] (390.0,512.0) rectangle (650.0,568.0);
\node[align=center,inner sep=0pt,outer sep=0pt,font=\fontsize{8.3}{11.62}\selectfont,text=monoLine] at (520.0,540.0) {一维管网水力};
\draw[draw=monoLine,line width=0.7pt,fill=white,rounded corners=0.1389cm] (390.0,578.0) rectangle (650.0,634.0);
\node[align=center,inner sep=0pt,outer sep=0pt,font=\fontsize{8.3}{11.62}\selectfont,text=monoLine] at (520.0,606.0) {二维地表漫流};
\draw[draw=monoLine,line width=0.7pt,fill=white,rounded corners=0.1389cm] (390.0,644.0) rectangle (650.0,700.0);
\node[align=center,inner sep=0pt,outer sep=0pt,font=\fontsize{8.3}{11.62}\selectfont,text=monoLine] at (520.0,672.0) {耦合求解调度};
\draw[draw=monoLine,line width=0.7pt,fill=white,rounded corners=0.1389cm] (390.0,710.0) rectangle (650.0,766.0);
\node[align=center,inner sep=0pt,outer sep=0pt,font=\fontsize{8.3}{11.62}\selectfont,text=monoLine] at (520.0,738.0) {残差智能订正};
\draw[line width=0.7pt,draw=monoLine] (330.0,573.0) -- (358.0,573.0);
\draw[line width=0.7pt,draw=monoLine] (358.0,540.0) -- (358.0,606.0);
\draw[line width=0.7pt,draw=monoLine] (358.0,540.0) -- (390.0,540.0);
\draw[line width=0.7pt,draw=monoLine] (358.0,606.0) -- (390.0,606.0);
\draw[line width=0.7pt,draw=monoLine] (330.0,705.0) -- (358.0,705.0);
\draw[line width=0.7pt,draw=monoLine] (358.0,672.0) -- (358.0,738.0);
\draw[line width=0.7pt,draw=monoLine] (358.0,672.0) -- (390.0,672.0);
\draw[line width=0.7pt,draw=monoLine] (358.0,738.0) -- (390.0,738.0);
\draw[line width=0.7pt,draw=monoLine] (650.0,540.0) -- (676.0,540.0);
\draw[line width=0.7pt,draw=monoLine] (650.0,606.0) -- (676.0,606.0);
\draw[line width=0.7pt,draw=monoLine] (650.0,672.0) -- (676.0,672.0);
\draw[line width=0.7pt,draw=monoLine] (650.0,738.0) -- (676.0,738.0);
\draw[line width=0.7pt,draw=monoLine] (676.0,540.0) -- (676.0,738.0);
\draw[line width=0.7pt,draw=monoLine] (676.0,639.0) -- (700.0,639.0);
\draw[draw=monoLine,line width=0.7pt,fill=white,rounded corners=0.1667cm] (700.0,512.0) rectangle (824.0,766.0);
\node[align=center,inner sep=0pt,outer sep=0pt,text=monoLine] at (762.0,639.0) {{\fontsize{8.80}{10.12}\selectfont\bfseries 内涝预报}\\[1.40pt]{\fontsize{7.00}{8.40}\selectfont 核心模型}};
\draw[fill=white,draw=monoLine,line width=0.7pt] (844.0,631.0) -- (864.0,631.0) -- (864.0,615.0) -- (884.0,639.0) -- (864.0,663.0) -- (864.0,647.0) -- (844.0,647.0) -- (844.0,663.0) -- (824.0,639.0) -- (844.0,615.0) -- cycle;
\draw[draw=monoLine,line width=0.7pt,fill=white,rounded corners=0.1667cm] (884.0,512.0) rectangle (946.0,766.0);
\node[align=center,inner sep=0pt,outer sep=0pt,font=\fontsize{8.3}{11.62}\selectfont,text=monoLine] at (915.0,592.8) {积};
\node[align=center,inner sep=0pt,outer sep=0pt,font=\fontsize{8.3}{11.62}\selectfont,text=monoLine] at (915.0,614.5) {水};
\node[align=center,inner sep=0pt,outer sep=0pt,font=\fontsize{8.3}{11.62}\selectfont,text=monoLine] at (915.0,636.2) {深};
\node[align=center,inner sep=0pt,outer sep=0pt,font=\fontsize{8.3}{11.62}\selectfont,text=monoLine] at (915.0,657.8) {度};
\node[align=center,inner sep=0pt,outer sep=0pt,font=\fontsize{8.3}{11.62}\selectfont,text=monoLine] at (915.0,679.5) {预};
\node[align=center,inner sep=0pt,outer sep=0pt,font=\fontsize{8.3}{11.62}\selectfont,text=monoLine] at (915.0,701.2) {警};
\draw[fill=white,draw=monoLine,line width=0.7pt] (546.0,404.0) -- (572.0,404.0) -- (572.0,446.0) -- (590.0,446.0) -- (559.0,482.0) -- (528.0,446.0) -- (546.0,446.0) -- cycle;
\draw[fill=white,draw=monoLine,line width=0.7pt] (546.0,764.0) -- (572.0,764.0) -- (572.0,806.0) -- (590.0,806.0) -- (559.0,842.0) -- (528.0,806.0) -- (546.0,806.0) -- cycle;
\draw[dashed,dash pattern=on 2.0pt off 1.6pt,line width=0.7pt,draw=monoLine] (40.0,842.0) rectangle (960.0,1086.0);
\draw[line width=0.7pt,draw=monoLine] (54.0,1072.0) -- (54.0,856.0) -- (86.0,856.0) -- (86.0,952.0);
\draw[line width=0.7pt,fill=white,draw=monoLine] (86.0,952.0) -- (118.0,952.0) -- (118.0,939.0) -- (160.0,964.0) -- (118.0,989.0) -- (118.0,976.0) -- (86.0,976.0);
\draw[line width=0.7pt,draw=monoLine] (86.0,976.0) -- (86.0,1072.0) -- (54.0,1072.0);
\node[align=center,inner sep=0pt,outer sep=0pt,font=\fontsize{8.6}{12.04}\selectfont,text=monoLine] at (70.0,895.0) {预};
\node[align=center,inner sep=0pt,outer sep=0pt,font=\fontsize{8.6}{12.04}\selectfont,text=monoLine] at (70.0,941.0) {警};
\node[align=center,inner sep=0pt,outer sep=0pt,font=\fontsize{8.6}{12.04}\selectfont,text=monoLine] at (70.0,987.0) {联};
\node[align=center,inner sep=0pt,outer sep=0pt,font=\fontsize{8.6}{12.04}\selectfont,text=monoLine] at (70.0,1033.0) {动};
\draw[fill=monoHead,draw=monoLine,line width=0.7pt,rounded corners=0.3472cm] (300.0,842.0) rectangle (700.0,892.0);
\node[align=center,inner sep=0pt,outer sep=0pt,font=\fontsize{8.8}{12.32}\selectfont,text=monoLine] at (500.0,867.0) {\bfseries 风险研判与分级预警};
\draw[draw=monoLine,line width=0.7pt,fill=white,rounded corners=0.1389cm] (300.0,906.0) rectangle (700.0,976.0);
\node[align=center,inner sep=0pt,outer sep=0pt,font=\fontsize{8.6}{12.04}\selectfont,text=monoLine] at (500.0,941.0) {\bfseries 内涝风险实时研判};
\draw[line width=0.7pt,draw=monoLine] (500.0,976.0) -- (500.0,1006.0);
\draw[line width=0.7pt,draw=monoLine] (288.0,1006.0) -- (828.0,1006.0);
\draw[line width=0.7pt,draw=monoLine] (288.0,1006.0) -- (288.0,1044.0);
\draw[draw=monoLine,line width=0.7pt,fill=white,rounded corners=0.1389cm] (208.0,1044.0) rectangle (368.0,1076.0);
\node[align=center,inner sep=0pt,outer sep=0pt,font=\fontsize{8.0}{11.20}\selectfont,text=monoLine] at (288.0,1060.0) {短信与 App 推送};
\draw[line width=0.7pt,draw=monoLine] (468.0,1006.0) -- (468.0,1044.0);
\draw[draw=monoLine,line width=0.7pt,fill=white,rounded corners=0.1389cm] (388.0,1044.0) rectangle (548.0,1076.0);
\node[align=center,inner sep=0pt,outer sep=0pt,font=\fontsize{8.0}{11.20}\selectfont,text=monoLine] at (468.0,1060.0) {泵站联合调度};
\draw[line width=0.7pt,draw=monoLine] (648.0,1006.0) -- (648.0,1044.0);
\draw[draw=monoLine,line width=0.7pt,fill=white,rounded corners=0.1389cm] (568.0,1044.0) rectangle (728.0,1076.0);
\node[align=center,inner sep=0pt,outer sep=0pt,font=\fontsize{8.0}{11.20}\selectfont,text=monoLine] at (648.0,1060.0) {交通诱导联动};
\draw[line width=0.7pt,draw=monoLine] (828.0,1006.0) -- (828.0,1044.0);
\draw[draw=monoLine,line width=0.7pt,fill=white,rounded corners=0.1389cm] (748.0,1044.0) rectangle (908.0,1076.0);
\node[align=center,inner sep=0pt,outer sep=0pt,font=\fontsize{8.0}{11.20}\selectfont,text=monoLine] at (828.0,1060.0) {抢险资源调度};
\draw[fill=white,draw=monoLine,line width=0.7pt] (546.0,1086.0) -- (572.0,1086.0) -- (572.0,1128.0) -- (590.0,1128.0) -- (559.0,1164.0) -- (528.0,1128.0) -- (546.0,1128.0) -- cycle;
\draw[draw=monoLine,line width=0.8pt,fill=monoHead,rounded corners=0.1111cm] (40.0,1164.0) rectangle (960.0,1232.0);
\node[align=center,inner sep=0pt,outer sep=0pt,font=\fontsize{9.2}{12.88}\selectfont,text=black] at (500.0,1198.0) {\bfseries 支撑「监测—推演—预警—处置」闭环的城市内涝实时防控};
\end{tikzpicture}
\end{center}

\newpage

\vspace{2pt}
{\footnotesize\textbf{(b) 彩色版}　同一套几何，仅替换配色：容器描边、圆柱、立体卡片、胶囊与结论条统一取主色系深浅两档。}

\vspace{0pt}
\begin{center}
\begin{tikzpicture}[x=0.0138889cm,y=-0.0138889cm]
\draw[dashed,dash pattern=on 2.0pt off 1.6pt,line width=0.7pt,draw=tealDeep] (40.0,40.0) rectangle (960.0,404.0);
\draw[line width=0.7pt,draw=tealDeep] (54.0,390.0) -- (54.0,54.0) -- (86.0,54.0) -- (86.0,210.0);
\draw[line width=0.7pt,fill=white,draw=tealDeep] (86.0,210.0) -- (118.0,210.0) -- (118.0,197.0) -- (160.0,222.0) -- (118.0,247.0) -- (118.0,234.0) -- (86.0,234.0);
\draw[line width=0.7pt,draw=tealDeep] (86.0,234.0) -- (86.0,390.0) -- (54.0,390.0);
\node[align=center,inner sep=0pt,outer sep=0pt,font=\fontsize{8.6}{12.04}\selectfont,text=tealDeep] at (70.0,153.0) {数};
\node[align=center,inner sep=0pt,outer sep=0pt,font=\fontsize{8.6}{12.04}\selectfont,text=tealDeep] at (70.0,199.0) {据};
\node[align=center,inner sep=0pt,outer sep=0pt,font=\fontsize{8.6}{12.04}\selectfont,text=tealDeep] at (70.0,245.0) {治};
\node[align=center,inner sep=0pt,outer sep=0pt,font=\fontsize{8.6}{12.04}\selectfont,text=tealDeep] at (70.0,291.0) {理};
\draw[fill=white,draw=tealDeep,line width=0.7pt] (168.0,68.0) -- (168.0,166.0) .. controls (168.0,171.0) and (346.0,171.0) .. (346.0,166.0) -- (346.0,68.0) -- cycle;
\draw[fill=white,draw=tealDeep,line width=0.7pt] (257.0,68.0) ellipse (89.0 and 9.0);
\node[align=center,inner sep=0pt,outer sep=0pt,font=\fontsize{8.3}{11.62}\selectfont,text=tealDeep] at (257.0,123.0) {雨量遥测库};
\draw[fill=white,draw=tealDeep,line width=0.7pt] (368.0,68.0) -- (368.0,166.0) .. controls (368.0,171.0) and (546.0,171.0) .. (546.0,166.0) -- (546.0,68.0) -- cycle;
\draw[fill=white,draw=tealDeep,line width=0.7pt] (457.0,68.0) ellipse (89.0 and 9.0);
\node[align=center,inner sep=0pt,outer sep=0pt,font=\fontsize{8.3}{11.62}\selectfont,text=tealDeep] at (457.0,123.0) {管网拓扑库};
\draw[fill=white,draw=tealDeep,line width=0.7pt] (568.0,68.0) -- (568.0,166.0) .. controls (568.0,171.0) and (746.0,171.0) .. (746.0,166.0) -- (746.0,68.0) -- cycle;
\draw[fill=white,draw=tealDeep,line width=0.7pt] (657.0,68.0) ellipse (89.0 and 9.0);
\node[align=center,inner sep=0pt,outer sep=0pt,font=\fontsize{8.3}{11.62}\selectfont,text=tealDeep] at (657.0,123.0) {历史内涝库};
\draw[fill=white,draw=tealDeep,line width=0.7pt] (768.0,68.0) -- (768.0,166.0) .. controls (768.0,171.0) and (946.0,171.0) .. (946.0,166.0) -- (946.0,68.0) -- cycle;
\draw[fill=white,draw=tealDeep,line width=0.7pt] (857.0,68.0) ellipse (89.0 and 9.0);
\node[align=center,inner sep=0pt,outer sep=0pt,font=\fontsize{8.3}{11.62}\selectfont,text=tealDeep] at (857.0,123.0) {预案知识库};
\draw[line width=0.7pt,draw=tealDeep] (257.0,196.0) -- (857.0,196.0);
\draw[line width=0.7pt,draw=tealDeep] (257.0,175.0) -- (257.0,196.0);
\draw[line width=0.7pt,draw=tealDeep] (457.0,175.0) -- (457.0,196.0);
\draw[line width=0.7pt,draw=tealDeep] (657.0,175.0) -- (657.0,196.0);
\draw[line width=0.7pt,draw=tealDeep] (857.0,175.0) -- (857.0,196.0);
\draw[line width=0.7pt,draw=tealDeep] (578.0,196.0) -- (578.0,222.0);
\fill[fill=white,draw=tealDeep,line width=0.7pt] (257.0,264.0) -- (299.0,222.0) -- (899.0,222.0) -- (857.0,264.0) -- cycle;
\fill[fill=tealSide,draw=tealDeep,line width=0.7pt] (857.0,264.0) -- (899.0,222.0) -- (899.0,332.0) -- (857.0,374.0) -- cycle;
\draw[fill=white,draw=tealDeep,line width=0.7pt] (257.0,264.0) rectangle (857.0,374.0);
\node[align=center,inner sep=0pt,outer sep=0pt,font=\fontsize{8.8}{12.32}\selectfont,text=tealDeep] at (578.0,243.0) {\bfseries 多源数据融合与样本构建};
\node[align=center,inner sep=0pt,outer sep=0pt,font=\fontsize{7.6}{10.64}\selectfont,text=black] at (557.0,310.0) {时钟对齐 · 缺测修复 · 汇水单元编码};
\node[align=center,inner sep=0pt,outer sep=0pt,font=\fontsize{7.6}{10.64}\selectfont,text=black] at (557.0,336.0) {特征构建 · 样本分层 · 训练验证切分};
\draw[dashed,dash pattern=on 2.0pt off 1.6pt,line width=0.7pt,draw=navyDeep] (40.0,482.0) rectangle (960.0,764.0);
\draw[line width=0.7pt,draw=navyDeep] (54.0,750.0) -- (54.0,496.0) -- (86.0,496.0) -- (86.0,611.0);
\draw[line width=0.7pt,fill=white,draw=navyDeep] (86.0,611.0) -- (118.0,611.0) -- (118.0,598.0) -- (160.0,623.0) -- (118.0,648.0) -- (118.0,635.0) -- (86.0,635.0);
\draw[line width=0.7pt,draw=navyDeep] (86.0,635.0) -- (86.0,750.0) -- (54.0,750.0);
\node[align=center,inner sep=0pt,outer sep=0pt,font=\fontsize{8.6}{12.04}\selectfont,text=navyDeep] at (70.0,554.0) {模};
\node[align=center,inner sep=0pt,outer sep=0pt,font=\fontsize{8.6}{12.04}\selectfont,text=navyDeep] at (70.0,600.0) {型};
\node[align=center,inner sep=0pt,outer sep=0pt,font=\fontsize{8.6}{12.04}\selectfont,text=navyDeep] at (70.0,646.0) {推};
\node[align=center,inner sep=0pt,outer sep=0pt,font=\fontsize{8.6}{12.04}\selectfont,text=navyDeep] at (70.0,692.0) {演};
\draw[draw=navyDeep,line width=0.7pt,fill=white,rounded corners=0.1667cm] (180.0,528.0) rectangle (330.0,618.0);
\node[align=center,inner sep=0pt,outer sep=0pt,text=navyDeep] at (255.0,573.0) {{\fontsize{8.30}{9.54}\selectfont\bfseries 模型输入}\\[1.20pt]{\fontsize{6.40}{7.68}\selectfont 边界与初始场}};
\draw[draw=navyDeep,line width=0.7pt,fill=white,rounded corners=0.1667cm] (180.0,660.0) rectangle (330.0,750.0);
\node[align=center,inner sep=0pt,outer sep=0pt,text=navyDeep] at (255.0,705.0) {{\fontsize{8.30}{9.54}\selectfont\bfseries 参数标定}\\[1.20pt]{\fontsize{6.40}{7.68}\selectfont 糙率 · 坡度 · 管径}};
\draw[line width=0.7pt,draw=navyDeep,-{Latex[length=1.5mm,width=1.0mm]}] (255,618.0) -- (255,660.0);
\draw[draw=navyDeep,line width=0.7pt,fill=white,rounded corners=0.1389cm] (390.0,512.0) rectangle (650.0,568.0);
\node[align=center,inner sep=0pt,outer sep=0pt,font=\fontsize{8.3}{11.62}\selectfont,text=navyDeep] at (520.0,540.0) {一维管网水力};
\draw[draw=navyDeep,line width=0.7pt,fill=white,rounded corners=0.1389cm] (390.0,578.0) rectangle (650.0,634.0);
\node[align=center,inner sep=0pt,outer sep=0pt,font=\fontsize{8.3}{11.62}\selectfont,text=navyDeep] at (520.0,606.0) {二维地表漫流};
\draw[draw=navyDeep,line width=0.7pt,fill=white,rounded corners=0.1389cm] (390.0,644.0) rectangle (650.0,700.0);
\node[align=center,inner sep=0pt,outer sep=0pt,font=\fontsize{8.3}{11.62}\selectfont,text=navyDeep] at (520.0,672.0) {耦合求解调度};
\draw[draw=navyDeep,line width=0.7pt,fill=white,rounded corners=0.1389cm] (390.0,710.0) rectangle (650.0,766.0);
\node[align=center,inner sep=0pt,outer sep=0pt,font=\fontsize{8.3}{11.62}\selectfont,text=navyDeep] at (520.0,738.0) {残差智能订正};
\draw[line width=0.7pt,draw=navyDeep] (330.0,573.0) -- (358.0,573.0);
\draw[line width=0.7pt,draw=navyDeep] (358.0,540.0) -- (358.0,606.0);
\draw[line width=0.7pt,draw=navyDeep] (358.0,540.0) -- (390.0,540.0);
\draw[line width=0.7pt,draw=navyDeep] (358.0,606.0) -- (390.0,606.0);
\draw[line width=0.7pt,draw=navyDeep] (330.0,705.0) -- (358.0,705.0);
\draw[line width=0.7pt,draw=navyDeep] (358.0,672.0) -- (358.0,738.0);
\draw[line width=0.7pt,draw=navyDeep] (358.0,672.0) -- (390.0,672.0);
\draw[line width=0.7pt,draw=navyDeep] (358.0,738.0) -- (390.0,738.0);
\draw[line width=0.7pt,draw=navyDeep] (650.0,540.0) -- (676.0,540.0);
\draw[line width=0.7pt,draw=navyDeep] (650.0,606.0) -- (676.0,606.0);
\draw[line width=0.7pt,draw=navyDeep] (650.0,672.0) -- (676.0,672.0);
\draw[line width=0.7pt,draw=navyDeep] (650.0,738.0) -- (676.0,738.0);
\draw[line width=0.7pt,draw=navyDeep] (676.0,540.0) -- (676.0,738.0);
\draw[line width=0.7pt,draw=navyDeep] (676.0,639.0) -- (700.0,639.0);
\draw[draw=navyDeep,line width=0.7pt,fill=white,rounded corners=0.1667cm] (700.0,512.0) rectangle (824.0,766.0);
\node[align=center,inner sep=0pt,outer sep=0pt,text=navyDeep] at (762.0,639.0) {{\fontsize{8.80}{10.12}\selectfont\bfseries 内涝预报}\\[1.40pt]{\fontsize{7.00}{8.40}\selectfont 核心模型}};
\draw[fill=white,draw=navyDeep,line width=0.7pt] (844.0,631.0) -- (864.0,631.0) -- (864.0,615.0) -- (884.0,639.0) -- (864.0,663.0) -- (864.0,647.0) -- (844.0,647.0) -- (844.0,663.0) -- (824.0,639.0) -- (844.0,615.0) -- cycle;
\draw[draw=navyDeep,line width=0.7pt,fill=white,rounded corners=0.1667cm] (884.0,512.0) rectangle (946.0,766.0);
\node[align=center,inner sep=0pt,outer sep=0pt,font=\fontsize{8.3}{11.62}\selectfont,text=navyDeep] at (915.0,592.8) {积};
\node[align=center,inner sep=0pt,outer sep=0pt,font=\fontsize{8.3}{11.62}\selectfont,text=navyDeep] at (915.0,614.5) {水};
\node[align=center,inner sep=0pt,outer sep=0pt,font=\fontsize{8.3}{11.62}\selectfont,text=navyDeep] at (915.0,636.2) {深};
\node[align=center,inner sep=0pt,outer sep=0pt,font=\fontsize{8.3}{11.62}\selectfont,text=navyDeep] at (915.0,657.8) {度};
\node[align=center,inner sep=0pt,outer sep=0pt,font=\fontsize{8.3}{11.62}\selectfont,text=navyDeep] at (915.0,679.5) {预};
\node[align=center,inner sep=0pt,outer sep=0pt,font=\fontsize{8.3}{11.62}\selectfont,text=navyDeep] at (915.0,701.2) {警};
\draw[fill=white,draw=navyDeep,line width=0.7pt] (546.0,404.0) -- (572.0,404.0) -- (572.0,446.0) -- (590.0,446.0) -- (559.0,482.0) -- (528.0,446.0) -- (546.0,446.0) -- cycle;
\draw[fill=white,draw=warmDeep,line width=0.7pt] (546.0,764.0) -- (572.0,764.0) -- (572.0,806.0) -- (590.0,806.0) -- (559.0,842.0) -- (528.0,806.0) -- (546.0,806.0) -- cycle;
\draw[dashed,dash pattern=on 2.0pt off 1.6pt,line width=0.7pt,draw=warmDeep] (40.0,842.0) rectangle (960.0,1086.0);
\draw[line width=0.7pt,draw=warmDeep] (54.0,1072.0) -- (54.0,856.0) -- (86.0,856.0) -- (86.0,952.0);
\draw[line width=0.7pt,fill=white,draw=warmDeep] (86.0,952.0) -- (118.0,952.0) -- (118.0,939.0) -- (160.0,964.0) -- (118.0,989.0) -- (118.0,976.0) -- (86.0,976.0);
\draw[line width=0.7pt,draw=warmDeep] (86.0,976.0) -- (86.0,1072.0) -- (54.0,1072.0);
\node[align=center,inner sep=0pt,outer sep=0pt,font=\fontsize{8.6}{12.04}\selectfont,text=warmDeep] at (70.0,895.0) {预};
\node[align=center,inner sep=0pt,outer sep=0pt,font=\fontsize{8.6}{12.04}\selectfont,text=warmDeep] at (70.0,941.0) {警};
\node[align=center,inner sep=0pt,outer sep=0pt,font=\fontsize{8.6}{12.04}\selectfont,text=warmDeep] at (70.0,987.0) {联};
\node[align=center,inner sep=0pt,outer sep=0pt,font=\fontsize{8.6}{12.04}\selectfont,text=warmDeep] at (70.0,1033.0) {动};
\draw[fill=warmLight,draw=warmDeep,line width=0.7pt,rounded corners=0.3472cm] (300.0,842.0) rectangle (700.0,892.0);
\node[align=center,inner sep=0pt,outer sep=0pt,font=\fontsize{8.8}{12.32}\selectfont,text=warmDeep] at (500.0,867.0) {\bfseries 风险研判与分级预警};
\draw[draw=warmDeep,line width=0.7pt,fill=white,rounded corners=0.1389cm] (300.0,906.0) rectangle (700.0,976.0);
\node[align=center,inner sep=0pt,outer sep=0pt,font=\fontsize{8.6}{12.04}\selectfont,text=warmDeep] at (500.0,941.0) {\bfseries 内涝风险实时研判};
\draw[line width=0.7pt,draw=warmDeep] (500.0,976.0) -- (500.0,1006.0);
\draw[line width=0.7pt,draw=warmDeep] (288.0,1006.0) -- (828.0,1006.0);
\draw[line width=0.7pt,draw=warmDeep] (288.0,1006.0) -- (288.0,1044.0);
\draw[draw=warmDeep,line width=0.7pt,fill=white,rounded corners=0.1389cm] (208.0,1044.0) rectangle (368.0,1076.0);
\node[align=center,inner sep=0pt,outer sep=0pt,font=\fontsize{8.0}{11.20}\selectfont,text=warmDeep] at (288.0,1060.0) {短信与 App 推送};
\draw[line width=0.7pt,draw=warmDeep] (468.0,1006.0) -- (468.0,1044.0);
\draw[draw=warmDeep,line width=0.7pt,fill=white,rounded corners=0.1389cm] (388.0,1044.0) rectangle (548.0,1076.0);
\node[align=center,inner sep=0pt,outer sep=0pt,font=\fontsize{8.0}{11.20}\selectfont,text=warmDeep] at (468.0,1060.0) {泵站联合调度};
\draw[line width=0.7pt,draw=warmDeep] (648.0,1006.0) -- (648.0,1044.0);
\draw[draw=warmDeep,line width=0.7pt,fill=white,rounded corners=0.1389cm] (568.0,1044.0) rectangle (728.0,1076.0);
\node[align=center,inner sep=0pt,outer sep=0pt,font=\fontsize{8.0}{11.20}\selectfont,text=warmDeep] at (648.0,1060.0) {交通诱导联动};
\draw[line width=0.7pt,draw=warmDeep] (828.0,1006.0) -- (828.0,1044.0);
\draw[draw=warmDeep,line width=0.7pt,fill=white,rounded corners=0.1389cm] (748.0,1044.0) rectangle (908.0,1076.0);
\node[align=center,inner sep=0pt,outer sep=0pt,font=\fontsize{8.0}{11.20}\selectfont,text=warmDeep] at (828.0,1060.0) {抢险资源调度};
\draw[fill=white,draw=warmDeep,line width=0.7pt] (546.0,1086.0) -- (572.0,1086.0) -- (572.0,1128.0) -- (590.0,1128.0) -- (559.0,1164.0) -- (528.0,1128.0) -- (546.0,1128.0) -- cycle;
\draw[draw=navyDeep,line width=0.8pt,fill=navyDeep,rounded corners=0.1111cm] (40.0,1164.0) rectangle (960.0,1232.0);
\node[align=center,inner sep=0pt,outer sep=0pt,font=\fontsize{9.2}{12.88}\selectfont,text=white] at (500.0,1198.0) {\bfseries 支撑「监测—推演—预警—处置」闭环的城市内涝实时防控};
\end{tikzpicture}
\end{center}

\vspace{0pt}
{\footnotesize
\textbf{图元}　数据库圆柱 ×4；立体折角卡片（顶面放标题、右侧面填灰）；胶囊标题；
分层虚线容器 + 左侧竖排标签 + 内嵌右向块箭头；层间空心块箭头；双向箭头；1→4 树形支架；结论条。
}

\end{document}
